  \documentclass[french,master,gcs]{webisthesis}

  \usepackage{booktabs}
  \usepackage{fontawesome}
  \usepackage{xcolor}
  \usepackage{soul}
  \usepackage[backend=biber,style=authoryear]{biblatex}

  %
  % Values
  % ------
  \ThesisSetTitle{Sécuriser l'autonomie : Cadre de défense des IA agentiques par le monitoring comportemental et les modèles souverains}
  \ThesisSetKeywords{These, are, my, Keywords} % only for PDF meta attributes
  \ThesisSetLocation{Lyon, France} 

  \ThesisSetAuthor{Clément RONDEAU}
  \ThesisSetStudentNumber{01234}
  \ThesisSetDateOfBirth{27}{4}{1992}
  \ThesisSetPlaceOfBirth{France}

  % Supervisors should usually be Professors from the candidate's university. A second supervisor is not always needed. 
  \ThesisSetSupervisors{Dr.}

  \ThesisSetSubmissionDate{30}{7}{2026}

  %
  % Utils
  % -------------------------------
  \newcommand{\bscom}[2]{
      \st{\scriptsize\,#1}{\color{blue}\scriptsize\,#2}%
    }
  % Create links in the pdf document
  \usehyperref

  \addbibresource{MasterThesisGuardia.bib}

  %
  % Document
  %

\begin{document}
\begin{frontmatter}
  \begin{abstract}
    Les IA agentiques se multiplient pour répondre à différents besoins au sein des entreprises. Aujourd'hui, on leur confie l'accès à des documents sensibles ainsi que la prise de décisions autonomes, ce qui en fait de nouvelles cibles critiques. Leur défense devient donc un sujet primordial pour limiter la surface d'attaque sur le système d'information. Ce mémoire explore la sécurisation de ces agents à travers le déploiement d'un monitoring comportemental en temps réel pour détecter les anomalies au runtime, telles que les injections de prompt. De plus, il démontre comment le choix de modèles souverains et frugaux, alignés sur la démarche "IA Frugale" de l'AFNOR, permet de réduire la dépendance aux API tierces, d'empêcher les fuites de données et de garantir un périmètre auditable. Cette approche est complétée par une preuve de concept (PoC) démontrant l'efficacité de ce cadre de durcissement.
  \end{abstract}

  \tableofcontents

  \chapter*{Remerciements} % optional
  Merci aux auteurs du modèle webisthesis pour leur travail.
  \\
  Je remercie également les auteurs des projets open-source LaTeX, Zotero, et l'ensemble des plugins utilisés pour la rédaction de ce document.

  % \listoffigures % optional, usually not needed

  % \listoftables % optional, usually not needed

  % \listofalgorithms % optional, usually not needed
  %    requires package algorithm2e

  % optional: list of symbols/notation (e.g., using the nomencl package) but usually not needed
\end{frontmatter}

\include{chapitre1}
% \include{chapter2}

% \include{appendixA}

% Bibliography
\nocite{*}
\printbibliography

\end{document}

